Seid ihr sicher, dass abdicate und abdication die korrekte Begriffe sind? Ich dachte, dass das heißt, dass eine Regierung oder ein Staatshaupt abtritt?
[sb] Ja, das ist tatsächlich in der Literatur der Fachbegriff dafür. Ich werde bei der ersten Verwendung des Begriffs noch mal eine Quelle hinzufügen, denn wenn Du darüber stolperst, dann auch andere. --> [update sb 29.08] Ich habe jetzt die konkrete Textstelle hinzugefügt (Bednar, S. 71f)


In Tabelle 2 & Tabelle 4 fehlen die Überschriften, die andeuten, was in jeder Spalte gezeigt wird. Diese findet man in den pdfs. Kann einer von euch die hinzufügen?
[sb] Habe sie eingefügt, das hatte ich leider vorher übersehen. 


Kommentar Salva: „%% I struggle with the table here. If the original variable is omitted and replaced with the 'relative difference'-variable, why is there still a number for the original variable in col (2)? SB 30.07.2024”
Antwort RvE: There should not be one, and I also do not see one. Maybe Julia took it out? Please check whether all is clear and OK now.
[sb] Ist geklärt und so passt es. 

Ich habe den folgenden Satz am Ende der Results umformuliert:
ALT: Prima facie, it cannot be ruled out that the results are biased by statements that were less intensely perceived, which could lead to misinterpretations of our findings.
NEU: Prima facie, it cannot be ruled out that prime ministers will have timed those statements that were meant to most clearly convey a message right around conferences, as they knew that statements at these moments would have reached the most attention and the largest possible audiences.
OK?
[sb] Super. Viel besser.

Salva: 3. Z. 169 / Beginn von Abschnitt 3: RvE: bitte überprüfen. Ich habe den Verweis auf Brand-width rausgenommen und statt dessen das twitter-api package zitiert, welches wir (anfäng-lich) tatsächlich benutzt haben. Julia: siehst Du hier Änderungsbedarf?
Nein, das passt aus meiner Sicht JR 24.08.2024
Ich kann nicht so gut abschätzen, in wieweit das so ausreicht. Sollten wir nicht noch etwas über die Keywords, die ihr benutzt habt sagen? RvE 27.08.2024
[sb] Ich habe eine Fußnote eingefügt, die das Verfahren beschreibt: We filtered out tweets related to COVID-19 on the accounts of the prime ministers or their press offices. OK Julia?


Julia: Z 251: Ich dachte die unabhängige Variable in Equation 1 ist die unabhängige Variable die Inzidenzrate (und nicht die Impfquote)?
[sb] Reyn?
RvE 27.08.2024: Habe ich nicht verstanden oder gefunden? (Vielleicht sind die Zeilennummer mittlerweile verschoben?)
[sb 28.08.] Julia, schaust Du noch mal, ob sich die Frage geklärt hat?

Salva: 16: Eine inhaltliche Frage: können oder sollten wir etwas Konkretes über den Effekt des Fe-deralism-Index sagen? Wenn ich es richtig verstehe, dann ist dieser Index als covariate ent-halten. Für die Conclusio wäre es aber schön, wir könnten etwas sagen wie: während die In-zidenz einen hochsignifikanten Effekt zeigt, finden wir keinen Zusammenhang zu den Grund-überzeugungen.

RvE 27.08.2024: I uploaded a table with the full results of the regressions in columns (4) and (5) of Table 1 (those are the ones that include this covariate). The coefficient is not significant in either regression.
If you want to write something about this, it should be along the lines of
… Appendix Table X shows the results from columns (4) and (5) of this table, which include the state-level controls. As can be seen from this table, the coefficients to none of the three variables reach significance. The coefficient to the variable measuring states’ general attitudes towards federalism suggests that states that had been strongly in favor of centralization before the pandemic became relatively less so during the pandemic, and vice versa.
[sb] Vielen Dank! Siehe Z. 399 und Appendix C.

Salva: 10. Z. 211 / Z. 322: Die Finanzkraftmesszahl ist in den state-level-controls berücksichtigt? Brauchen wir sie überhaupt? Oder sollten wir darauf verweisen, dass die Finanzkraft die 'ge-neral stances towards federalism' erklärt. Kurz: sollten wir den Teil mit der Finanzkraft kür-zen/streichen oder so belassen?
Ich denke das erste was sich eine, insbesondere eine ökonomisch versierte Leserschaft fragt, ist, inwiefern die Finanzstärke der Länder eine Rolle spielt. Daher würde ich zumindest in ei-nem Nebensatz erwähnen, dass die „general stance towards federalism“ auch die Finanzkraft berücksichtigt. JR 24.08.24
RvE 27.08.24: Wenn man die Variable aus der Regression weglässt, dann ändert sich der geschätzte Effekt nicht (wird ganz leicht größer: 2.01 statt 1.98). Oder meintest du das nicht? Sonst ist das eine reine inhaltliche Entscheidung.
[sb] Das ist auf jeden Fall beruhigend. Ich denke noch mal darüber nach (Stand: 28.08.)
